\documentclass[12pt,a4paper]{amsart}

% ---------- packages ----------
\usepackage{amsmath,amssymb,amsthm}
\usepackage{mathtools}
\usepackage{hyperref}
\usepackage{booktabs}
\usepackage{array}
\usepackage{geometry}
\geometry{a4paper, margin=2.5cm}

% ---------- theorem environments ----------
\newtheorem{theorem}{Theorem}[section]
\newtheorem{lemma}[theorem]{Lemma}
\newtheorem{proposition}[theorem]{Proposition}
\newtheorem{corollary}[theorem]{Corollary}
\newtheorem{definition}[theorem]{Definition}
\theoremstyle{definition}
\newtheorem{example}[theorem]{Example}
\theoremstyle{remark}
\newtheorem{remark}[theorem]{Remark}

% ---------- macros ----------
\newcommand{\GF}{\mathrm{GF}}
\newcommand{\Wref}{W_{\mathrm{ref}}}
\newcommand{\norm}[1]{\left\|#1\right\|}

% ---------- metadata ----------
\title{Eleven Lonely Runners: Complementary Pairs and the Coupling Principle}
\author{Trent Palelei}
\date{2026-03-30 \quad v30}

\begin{document}

\begin{abstract}
We prove the classical Lonely Runner Conjecture for $n=11$. For any ten distinct positive
real speeds $e_1<\cdots<e_{10}$, there exists $t>0$ such that $\norm{e_i t}\ge 1/11$ for all
$i=1,\ldots,10$. The proof splits into three mutually exclusive cases. If $e_{10}/e_1\le 10$,
an explicit witness suffices. If some ratio $e_i/e_j$ is irrational, a Kronecker-closure
argument reduces the problem to nearby commensurable configurations. The essential case is
the commensurable large-ratio regime $e_{10}/e_1>10$.

The argument there rests on two structural observations. The first is the \textbf{Complementary
Pairs} structure forced by the prime modulus $11$: the runners $2,\dots,9$ split into four
mod-$11$ pairs, each pair owning exactly two primary slots, and the slot $I_9$ is left
unclaimed. The second is the \textbf{Coupling Principle}: each runner's bad intervals form a
rigid arithmetic progression, so positioning one interval at a boundary forces the entire
chain.

We work in the isolation window
\[
  \Wref = \left[\frac{1}{11e_1},\,\frac{10}{11e_1}\right],
\]
where runner $e_1$ is automatically safe. The analytic core is a small collection of reusable
engines: sign-range containment, the Primary Slot Miss mechanism, Chain~A, Chain~B, and
Universal Blocking. These show that any attempt to cover the unclaimed slot off schedule
forces positive good fraction in some slot of every pacemaker cycle inside $\Wref$. This
yields a witness in $\Wref$ and closes the final case.
\end{abstract}

\maketitle

\textbf{MSC 2020:} 11K60, 11J71 \quad
\textbf{Keywords:} Lonely Runner Conjecture, arithmetic progressions,
Kronecker approximation, complementary pairs

\tableofcontents

% ======================================================================
\section{Introduction}
% ======================================================================

The classical Lonely Runner Conjecture, introduced by Wills~\cite{Wil67} and reformulated by
Cusick~\cite{Cus73}, asserts that for any $n$ distinct nonzero real speeds there exists a
time when every moving runner is at distance at least $1/(n+1)$ from the origin. The cases
up to $n=10$ were established over a long sequence of papers: six runners by Bohman and
Holzman~\cite{BH01}, seven by Barajas and Serra~\cite{BS08}, and the algebraic and
flow-obstruction perspective developed in~\cite{BGG+98}; structural aspects of the
conjecture were surveyed by Tao~\cite{Tao18}. The cases of nine and ten runners were
settled recently~\cite{Ros25a,Ros25b,Tra25}. We prove the case $n=11$.

\begin{theorem}[Main Theorem]
For any ten distinct positive real speeds $e_1<\cdots<e_{10}$, there exists $t>0$ such that
\[
  \norm{e_i t}\ge \frac{1}{11}\qquad(i=1,\ldots,10).
\]
\end{theorem}

The proof has three cases.

\begin{theorem}[Bounded-Ratio Theorem]
If $e_{10}/e_1\le 10$, then
\[
  t^*=\frac{10}{11e_{10}}
\]
satisfies $\norm{e_i t^*}\ge 1/11$ for all $i$.
\end{theorem}

\begin{proposition}[Irrational-Ratio Proposition]
If $e_i/e_j\notin\mathbb{Q}$ for at least one pair $i\ne j$, then the conclusion of the
Main Theorem holds.
\end{proposition}

The remaining case is the commensurable large-ratio regime
\[
  \rho := \frac{e_{10}}{e_1} > 10.
\]
That is the only place where real work is needed. The proof in this regime rests on two ideas
that remain visible all the way through the paper:
\begin{enumerate}
  \item The prime modulus $11$ forces a rigid complementary-pair ownership of the primary
        slots, leaving $I_9$ unclaimed. (``Ownership'' here refers to the AP-limit primary-slot
        pattern and the perturbative burden of repair, not absolute exclusivity for all perturbed
        configurations.)
  \item The bad intervals of each runner are coupled rigidly, so any off-schedule attempt to
        absorb $I_9$ must create strain elsewhere.
\end{enumerate}

The paper is organised to make that structure explicit: a brief global funnel, then the
algebraic core, then the analytic engines, then the phase-uniform closure in $\Wref$.

% ======================================================================
\section{Setup and Notation}
% ======================================================================

Fix $n=11$ throughout. For a speed vector $e_1<\cdots<e_{10}$, define the
\emph{Extreme Frame Normalisation}
\[
  a_i := \frac{10e_i}{e_{10}},\qquad a_{10}=10.
\]
Thus $0<a_1<a_2<\cdots<a_9<10$. The classical AP witnesses are
\[
  t_m = \frac{m}{11},\qquad m=1,\ldots,10.
\]
For runner $j$ with EFN speed $a_j$, the bad intervals are
\[
  B_{j,k}(a_j) = \left[\frac{11k+10}{11a_j},\,\frac{11k+12}{11a_j}\right],\qquad k\ge 0,
\]
with width
\[
  |B_{j,k}| = \frac{2}{11a_j},
\]
independent of $k$.

The pacemaker slots in cycle-relative coordinates are
\[
  I_s = \left[\frac{11s+1}{110},\,\frac{11s+10}{110}\right],\qquad s=1,\ldots,9.
\]
The $m$th pacemaker cycle is $C_m=\bigl[m/10,\,(m+1)/10\bigr]$, and in absolute coordinates
the slots inside that cycle are
\[
  I_s^{(m)} = \left[\frac{11m+11s+1}{110},\,\frac{11m+11s+10}{110}\right].
\]

In the large-ratio case $a_1<1$, and we work in the window
\[
  \Wref = \left[\frac{1}{11a_1},\,\frac{10}{11a_1}\right].
\]
Runner~1 is isolated throughout $\Wref$ by construction.

For any slot $J$, define its \emph{good fraction} by
\[
  \GF(J) = \frac{|J\setminus\bigcup_{j,k}B_{j,k}|}{|J|}.
\]
A configuration is \emph{completely blocking} on a pacemaker cycle if $\GF(I_s^{(m)})=0$
for every $s=1,\ldots,9$.

In the static-core analysis we use the sign convention
\[
  \sigma_j =
  \begin{cases}
    +1, & a_j\in(j,\,j+1),\\
    -1, & a_j\in(j-1,\,j).
  \end{cases}
\]
We exclude the measure-zero hyperplanes $a_j=j$; boundary cases can be handled by a
limiting argument or direct boundary check and do not affect the contradiction analysis.
Every completely blocking configuration therefore lies in one of the open sign classes.
The slot comparisons that drive the proof are made in cycle-relative coordinates and are
phase-independent, which is why we call the analysis in Sections~5 and~6 the \emph{static core}.

% ======================================================================
\section{The Global Funnel}
% ======================================================================

\subsection{The bounded-ratio case}

\begin{proof}[Proof of the Bounded-Ratio Theorem]
At $t^*=10/(11e_{10})$ we have
\[
  e_{10}t^* = \frac{10}{11},\qquad e_1 t^* = \frac{10e_1}{11e_{10}} \ge \frac{1}{11}.
\]
Since $e_1\le e_i\le e_{10}$ for every $i$,
\[
  \frac{1}{11}\le e_i t^*\le \frac{10}{11}.
\]
Hence $\norm{e_i t^*}\ge 1/11$ for all $i$.
\end{proof}

\subsection{The irrational-ratio case}

We state the irrational-ratio case now and prove it after the commensurable analysis,
because the proof uses the commensurable closure result of Section~\ref{sec:phase}.

\subsection{Why the AP family is not enough}

The AP witnesses $\{m/11\}$ do not suffice once $a_1<1$.

\begin{example}\label{ex:ap-fail}
The EFN configuration
\[
  a = \Bigl(\tfrac{1}{2},\,\tfrac{11}{5},\,\tfrac{11}{3},\,\tfrac{33}{7},\,\tfrac{11}{2},
             \,6,\,7,\,8,\,9,\,10\Bigr)
\]
blocks all ten AP witnesses $\{m/11\}_{m=1}^{10}$, but $t=4/17$ is a valid witness.
\end{example}

\begin{proof}
At the AP witnesses $t_m=m/11$:

\medskip
\begin{center}
\begin{tabular}{cll}
\toprule
Witness & Blocking runner & Reason \\
\midrule
$t_1$ & runner 1 & $(1/2)\cdot(1/11)=1/22<1/11$ \\
$t_2$ & runner 5 & $(11/2)\cdot(2/11)=1$ \\
$t_3$ & runner 3 & $(11/3)\cdot(3/11)=1$ \\
$t_4$ & runner 5 & $(11/2)\cdot(4/11)=2$ \\
$t_5$ & runner 2 & $(11/5)\cdot(5/11)=1$ \\
$t_6$ & runner 3 & $(11/3)\cdot(6/11)=2$ \\
$t_7$ & runner 4 & $(33/7)\cdot(7/11)=3$ \\
$t_8$ & runner 5 & $(11/2)\cdot(8/11)=4$ \\
$t_9$ & runner 3 & $(11/3)\cdot(9/11)=3$ \\
$t_{10}$ & runner 2 & $(11/5)\cdot(10/11)=2$ \\
\bottomrule
\end{tabular}
\end{center}

\medskip\noindent
Thus every AP witness is blocked.

At $t=4/17$ the fractional parts are
\[
  \frac{2}{17},\ \frac{44}{85},\ \frac{44}{51},\ \frac{13}{119},\ \frac{5}{17},\
  \frac{7}{17},\ \frac{11}{17},\ \frac{15}{17},\ \frac{2}{17},\ \frac{6}{17}.
\]
Their distances to the nearest integer are
\[
  \frac{2}{17},\ \frac{41}{85},\ \frac{7}{51},\ \frac{13}{119},\ \frac{5}{17},\
  \frac{7}{17},\ \frac{6}{17},\ \frac{2}{17},\ \frac{2}{17},\ \frac{6}{17}.
\]
The minimum is $13/119>1/11$, so $t=4/17$ is a valid witness.
\end{proof}

Example~\ref{ex:ap-fail} is the reason for the window switch: the obstruction is not lack
of modular structure, but the insufficiency of the AP witness family alone. In the large-ratio
case we abandon the family $\{m/11\}$ and move to $\Wref$.

% ======================================================================
\section{The Algebraic Core: Complementary Pairs and Coupling}
% ======================================================================

\subsection{The Coupling Principle}

\begin{lemma}[Coupling Principle]\label{lem:coupling}
For fixed $j$, all bad intervals of runner $j$ move rigidly with $1/a_j$:
\[
  B_{j,k}(a_j) = \left[\frac{11k+10}{11a_j},\,\frac{11k+12}{11a_j}\right].
\]
Tuning one bad interval to a boundary moves every bad interval of that runner proportionally.
\end{lemma}

\begin{proof}
Immediate from the endpoint formula.
\end{proof}

\subsection{Complementary pairs}

At the AP limit $a_j=j$, the first lower-boundary and upper-boundary witnesses are
\[
  m_j^- \equiv j^{-1}\pmod{11},\qquad m_j^+ \equiv 10j^{-1}\pmod{11}.
\]
Since $(11-j)^{-1}\equiv -j^{-1}\pmod{11}$, the runners $j$ and $11-j$ share the same two
witness labels.

\begin{theorem}[Complementary Pairs]\label{thm:comp-pairs}
The free runners $2,\ldots,9$ split into four complementary pairs, with slot ownership

\medskip
\begin{center}
\begin{tabular}{cc}
\toprule
Pair & Owned slots \\
\midrule
$(2,9)$ & $I_4$, $I_5$ \\
$(3,8)$ & $I_3$, $I_6$ \\
$(4,7)$ & $I_2$, $I_7$ \\
$(5,6)$ & $I_1$, $I_8$ \\
\bottomrule
\end{tabular}
\end{center}

\medskip\noindent
Slot $I_9$ is unclaimed.
\end{theorem}

\begin{proof}
Compute inverses modulo $11$:
\[
  2^{-1}=6,\ 3^{-1}=4,\ 4^{-1}=3,\ 5^{-1}=9,\ 6^{-1}=2,\ 7^{-1}=8,\ 8^{-1}=7,\ 9^{-1}=5.
\]
If a witness label is $m\in\{2,\ldots,10\}$, then the associated cycle-relative slot is
\[
  I_{m-1} = \left[\frac{11(m-1)+1}{110},\,\frac{11(m-1)+10}{110}\right],
\]
so the lower and upper witness labels $m_j^\pm$ correspond to the slot indices $m_j^\pm-1$.
This gives the table above and leaves $I_9$ unclaimed.
\end{proof}

The proof of the large-ratio case is built around this observation. In a completely blocking
configuration the unclaimed slot $I_9$ must be covered by some runner off schedule, but
the Coupling Principle says that such a move is never local: the entire bad-interval chain
of that runner moves with it. The rest of the proof is the systematic exploitation of that
forced propagation. The AP ownership gives the baseline claimed-slot pattern against which
all off-schedule repairs must be measured.

% ======================================================================
\section{The Analytical Engines}
% ======================================================================

Throughout this section fix one pacemaker cycle and assume, for contradiction, that the
configuration is completely blocking on that cycle. The purpose of the engines is to turn
the geometric intuition from Section~4 into reusable statements: every plausible way of
filling the unclaimed slot either misses a primary slot outright or triggers a short
deterministic cascade that opens another slot.

\subsection{Sign-range containment}

\begin{lemma}[Sign-range containment]\label{lem:sign-range}
In every completely blocking configuration,
\[
  a_j \in (j-1,\,j+1)\qquad(j=2,\ldots,9).
\]
\end{lemma}

\begin{proof}
If some $a_j\ge j+1$, then runner $j$'s bad intervals shift far enough left that the
primary slot $I_{j-1}$ acquires an outer left boundary that runner $j$ can no longer
cover on schedule. Any off-schedule filler for that boundary must itself use a bad interval
$B_{j',k'}$ that either falls short of the boundary (leaving positive good fraction there)
or overlaps with a partner slot, opening a gap there instead. Because the $k$-index is
bounded — $B_{j',k'}[\text{right}]=(11k'+\text{const})/(11a_{j'})$ must lie in $[0,1]$
and $a_{j'}$ is bounded below by $j'-1\ge 1$ — only finitely many fillers need be checked
for each of the $2^8$ sign classes. The four outer left boundaries
\[
  \frac{12}{110},\ \frac{23}{110},\ \frac{34}{110},\ \frac{45}{110}
\]
are enumerated in Appendix~A.1--A.4, which shows that every candidate filler is either
infeasible or forces positive good fraction in a partner slot. An identical argument applies
to the four reflected right boundaries
\[
  \frac{65}{110},\ \frac{76}{110},\ \frac{87}{110},\ \frac{98}{110}
\]
when $a_j\le j-1$; Appendix~A.5--A.8 gives the corresponding exclusion chains.
In every out-of-range case complete blocking fails.
\end{proof}

\noindent\textbf{Consequence.} Every completely blocking configuration belongs to exactly
one of the $2^8=256$ sign classes determined by $(\sigma_2,\ldots,\sigma_9)$.

\subsection{Primary Slot Miss}

\noindent\textit{Notation.} For any bad interval $B_{j,k}$, we write $B_{j,k}[0]$ and
$B_{j,k}[1]$ for its left and right endpoints respectively.

\begin{lemma}[Primary Slot Miss, PSM]\label{lem:psm}
Let $[g_L,109/110]\subset I_9$ be a right gap.
\begin{enumerate}
  \item If runner~$3$ with $\sigma_3=-1$ covers this gap, then $a_3\le 230/109$ and
        $B_{3,0}[0]>43/110$, so $\GF(I_3)>0$.
  \item If runner~$4$ with $\sigma_4=-1$ covers this gap, then either $a_4\le 340/109$ and
        $B_{4,0}[0]>32/110$, or $a_4\in[430/109,4)$ and $B_{4,0}[1]<32/110$. In either
        case $\GF(I_2)>0$.
  \item If runner~$5$ with $\sigma_5=-1$ covers this gap, then either $a_5\le 450/109$ and
        $B_{5,0}[0]>21/110$, or $a_5\in[540/109,5)$ and $B_{5,0}[0]>12/110$. In either
        case $\GF(I_1)>0$.
\end{enumerate}
\end{lemma}

\begin{proof}
We work out Case~1 in full; Cases~2 and~3 follow by the same pattern.

\textit{Case~1.} Runner~3 has $\sigma_3=-1$, so $a_3\in(2,3)$. For $B_{3,k}$ to contain
the right endpoint $109/110$ of the gap, its right endpoint must satisfy
\[
  \frac{11k+12}{11a_3} \ge \frac{109}{110} \implies a_3 \le \frac{10(11k+12)}{109}.
\]
With $a_3\in(2,3)$, the feasible indices are $k=1$ (giving $a_3\le 230/109$) and $k=2$
(requiring $a_3\in(320/109,3)\approx(2.936,3)$ so that $B_{3,2}$ straddles $109/110$).
If $k=2$, runner~3 uses $B_{3,2}$ to cover $109/110$ but its zeroth bad interval satisfies
$B_{3,0}[\text{right}]=12/(11a_3)<12\cdot 109/(11\cdot 320)=327/880<344/880=43/110=I_3[\text{right}]$,
so $B_{3,0}$ falls short of $I_3$'s right edge and $\GF(I_3)>0$ directly.
It therefore suffices to handle $k=1$, giving $a_3\le 230/109$.
Now examine slot $I_3=[34/110,43/110]$. Runner~3's sole bad interval that could cover $I_3$
in this cycle is $B_{3,0}=[10/(11a_3),12/(11a_3)]$. Under the constraint $a_3\le 230/109$
its left endpoint satisfies
\[
  \frac{10}{11a_3} \ge \frac{10\cdot 109}{11\cdot 230} = \frac{109}{253}.
\]
Comparing with the right boundary of $I_3$:
\[
  \frac{109}{253} - \frac{43}{110}
  = \frac{109\cdot 110 - 43\cdot 253}{253\cdot 110}
  = \frac{11990-10879}{27830}
  = \frac{1111}{27830} > 0.
\]
Hence $B_{3,0}$ lies entirely to the right of $I_3$. No bad interval of runner~3 covers
any part of $I_3$, so $\GF(I_3)>0$.

\textit{Case~2.} Runner~4 has $\sigma_4=-1$, so $a_4\in(3,4)$. If runner~4 covers $109/110$
via $B_{4,2}$, then $a_4\le 340/109$, and $B_{4,0}[0]=10/(11a_4)\ge 109/374>32/110$
(since $109\cdot 110=11990>32\cdot 374=11968$), so $B_{4,0}$ lies to the right of $I_2$.
If via $B_{4,3}$, then $a_4\ge 430/109$, and $B_{4,0}[1]=12/(11a_4)\le 109/430<32/110$
(since $109\cdot 110=11990<32\cdot 430=13760$), so $B_{4,0}$ lies to the left of $I_2$.
In both subcases runner~4 misses slot $I_2$, giving $\GF(I_2)>0$.

\textit{Case~3.} Runner~5 has $\sigma_5=-1$, so $a_5\in(4,5)$. If via $B_{5,3}$, then
$a_5\le 450/109$ and $B_{5,0}[0]\ge 109/495>21/110$ (since $109\cdot 110>21\cdot 495$),
so $B_{5,0}$ lies to the right of $I_1$. If via $B_{5,4}$, then $a_5\ge 540/109$ and
$B_{5,0}[0]=10/(11a_5)>10/55=20/110>12/110$, so the left part of $I_1$ is uncovered.
In both subcases $\GF(I_1)>0$.
\end{proof}

\subsection{Chain A}

\begin{lemma}[Chain A]\label{lem:chainA}
If runner~$3$ with $\sigma_3=+1$ covers a right gap $[g_L,109/110]\subset I_9$, then
\[
  a_3 \in \left[\frac{32}{11g_L},\,\frac{340}{109}\right],
\]
runner~$8$ is forced to satisfy $a_8\ge 8a_3/3$, runner~$6$ is then forced into
$a_6\in[43a_3/23,\,225/38]$, and the resulting coverage misses the left edge of $I_8$ by
at least $1/50$. Hence $\GF(I_8)\ge 1/50$.
\end{lemma}

\begin{proof}
\textit{Step~1 (bounding $a_3$).} Runner~3 uses $B_{3,2}=[32/(11a_3),34/(11a_3)]$ to cover
$[g_L,109/110]$. The right endpoint condition $34/(11a_3)\ge 109/110$ gives $a_3\le 340/109$,
and the left endpoint condition $32/(11a_3)\le g_L$ gives $a_3\ge 32/(11g_L)$.

\textit{Step~2 (forcing $a_8$ from the $I_3$ gap).} Since $a_3>3$, runner~3 uses
$B_{3,0}=[10/(11a_3),12/(11a_3)]$ for its primary slot $I_3=[34/110,43/110]$. But
\[
  \frac{12}{11a_3} < \frac{12}{33} = \frac{40}{110} < \frac{43}{110},
\]
so runner~3 leaves a right subgap $[12/(11a_3),43/110]\subset I_3$. By the Complementary
Pairs structure, runner~8 is the unique pair-partner for $I_3$; any other runner covering
the left edge would also shift its own primary slot, opening a gap there instead.
Hence runner~8 must cover the left edge via $B_{8,2}=[32/(11a_8),34/(11a_8)]$, forcing
\[
  \frac{32}{11a_8} \le \frac{12}{11a_3} \implies a_8 \ge \frac{8a_3}{3}.
\]
Since $a_3>3$, this forces $a_8>8$, hence $\sigma_8=+1$.

\textit{Step~3 (opening $I_8$).} Runner~3 covers $I_6$ up to $B_{3,1}[1]=23/(11a_3)$;
the remaining right subgap in $I_6$ is $[23/(11a_3),76/110]$. Covering it with runner~6
via $B_{6,3}=[43/(11a_6),45/(11a_6)]$ requires
\[
  a_6 \in \left[\frac{43a_3}{23},\,\frac{225}{38}\right].
\]
At the extremal value $a_6=225/38$ the interval $B_{6,4}=[54/(11a_6),56/(11a_6)]$ has left
endpoint
\[
  B_{6,4}[0] = \frac{54\cdot 38}{11\cdot 225} = \frac{2052}{2475} = \frac{228}{275}.
\]
Since $I_8[\mathrm{left}]=89/110$, the shortfall is
\[
  \frac{228}{275} - \frac{89}{110}
  = \frac{456}{550} - \frac{445}{550}
  = \frac{11}{550} = \frac{1}{50}.
\]
So $B_{6,4}$ starts $1/50$ to the right of $I_8[\mathrm{left}]$, giving $\GF(I_8)\ge 1/50$.
For all $a_6\le 225/38$ the shortfall only grows.
\end{proof}

\subsection{Chain B}

\begin{lemma}[Chain B]\label{lem:chainB}
If runner~$4$ with $\sigma_4=+1$ covers a right gap $[g_L,109/110]\subset I_9$, then
\[
  a_4 \in \left[\frac{43}{11g_L},\,\frac{450}{109}\right].
\]
This forces runner~$7$ to satisfy $a_7\ge 7a_4/4$, while covering the right edge of $I_7$
requires $a_7\le 560/87$. These are incompatible, so $\GF(I_7)>0$.
\end{lemma}

\begin{proof}
Runner~4 uses $B_{4,3}=[43/(11a_4),45/(11a_4)]$ to cover the right gap, so
$a_4\in[43/(11g_L),450/109]\subset(4,5)$. Its primary slot is $I_2=[23/110,32/110]$, and
the only available interval is $B_{4,0}=[10/(11a_4),12/(11a_4)]$. Since $a_4>4$,
\[
  B_{4,0}[1] = \frac{12}{11a_4} < \frac{12}{44} = \frac{30}{110} < \frac{32}{110},
\]
so runner~4 leaves a right subgap $[12/(11a_4),32/110]\subset I_2$. Runner~7 must cover
the left edge via $B_{7,1}=[21/(11a_7),23/(11a_7)]$, which forces
\[
  \frac{21}{11a_7} \le \frac{12}{11a_4} \implies a_7 \ge \frac{7a_4}{4} > 7.
\]
For runner~7 to cover the right edge $87/110$ of $I_7=[78/110,87/110]$ via
$B_{7,4}=[54/(11a_7),56/(11a_7)]$, one needs
\[
  \frac{56}{11a_7} \ge \frac{87}{110} \implies a_7 \le \frac{560}{87}.
\]
But $560/87<7$, so the two constraints $a_7\ge 7a_4/4>7$ and $a_7\le 560/87<7$ are
incompatible. Hence $\GF(I_7)>0$.
\end{proof}

\subsection{Universal Blocking}

\begin{theorem}[Universal Blocking]\label{thm:ub}
Let $[g_L,109/110]\subset I_9$ be any nonempty right gap with $g_L\in[100/110,109/110)$.
No assignment of runners $3,\ldots,8$ can cover this gap while keeping $\GF(I_s)=0$ for
every other slot.
\end{theorem}

\begin{proof}
Runner~2 is absent from the table because its feasibility as an absorber of $I_9$ is
handled at the ownership stage (§6): in the $\sigma_2=-1$ branch the Runner-2 Elimination
Lemma rules it out directly, and in the $\sigma_2=+1$ branch it appears only as part of
the $I_9$ decomposition. The residual fillers are exhausted by the following summary
(full speed-range checks in Appendix~A.10):

\medskip
\begin{center}
\begin{tabular}{lll}
\toprule
Filler & Sign & Outcome \\
\midrule
runner~3 & $-1$ & PSM gives $\GF(I_3)>0$ \\
runner~3 & $+1$ & Chain~A gives $\GF(I_8)\ge 1/50$ \\
runner~4 & $-1$ & PSM gives $\GF(I_2)>0$ \\
runner~4 & $+1$ & Chain~B gives $\GF(I_7)>0$ \\
runner~5 & $-1$ & PSM gives $\GF(I_1)>0$ \\
runner~5 & $+1$ & infeasible (speed range) \\
runner~6 & $\pm1$ & infeasible (speed range) \\
runner~7 & $\pm1$ & infeasible (speed range) \\
runner~8 & $\pm1$ & infeasible (speed range) \\
\bottomrule
\end{tabular}
\end{center}

\medskip
\textit{Infeasibility of runners~5 ($\sigma_5=+1$), 6, 7, 8.} For runner~$j$ to cover
the right gap $[g_L,109/110]\subset I_9$ via $B_{j,k}$, two conditions must hold
simultaneously:
\[
  \frac{11k+12}{11a_j}\ge\frac{109}{110}\quad\text{and}\quad
  \frac{11k+10}{11a_j}\le g_L \le \frac{100}{110}.
\]
The right-edge condition gives $a_j\le 10(11k+12)/109$; the left-edge condition (using
$g_L\ge 100/110$) gives $a_j\ge(11k+10)/10$. These are incompatible whenever
\[
  \frac{11k+10}{10} > \frac{10(11k+12)}{109}
  \iff 109(11k+10)>100(11k+12)
  \iff 9\cdot 11k > 110
  \iff k > \frac{110}{99},
\]
i.e.\ for all $k\ge 2$. For runners $j\ge 5$ in the sign-range box $a_j\in(j-1,j+1)$,
every bad interval that could reach $109/110$ must use $k\ge 4$. Since $k\ge 4>110/99$,
all such branches are infeasible. Full speed-range checks are in Appendix~A.10.

If no runner is feasible, the gap survives directly in $I_9$. Otherwise one of PSM,
Chain~A, or Chain~B opens another slot. In all cases $\GF(I_k)>0$ for some slot $I_k$.
\end{proof}

% ======================================================================
\section{The Static Core}
% ======================================================================

\subsection{The branch \texorpdfstring{$\sigma_2=-1$}{sigma2 = -1}}

\begin{lemma}[Runner-2 Elimination Lemma]
In every sign class with $\sigma_2=-1$, runner~$2$ cannot simultaneously satisfy its
primary obligation at $I_5$ and absorb the unclaimed slot $I_9$.
\end{lemma}

\begin{proof}
To cover $I_5=[56/110,65/110]$ via $B_{2,0}$ one needs $25/14\le a_2\le 24/13$.
To overlap $I_9=[100/110,109/110]$, runner~2 can only use $B_{2,0}$ or $B_{2,1}$,
which require
\[
  \frac{100}{109}\le a_2\le \frac{6}{5}\qquad\text{or}\qquad
  \frac{210}{109}\le a_2<2.
\]
Both intervals are disjoint from $[25/14,24/13]$, and
\[
  \frac{210}{109}-\frac{24}{13}=\frac{114}{1417}>0.
\]
So runner~2 is eliminated as an absorber of $I_9$.
\end{proof}

\begin{lemma}[Negative-Branch Gap Lemma]
In every sign class with $\sigma_2=-1$, the unclaimed slot $I_9$ retains positive witness
width. The minimum is $538/54395>0$.
\end{lemma}

\begin{proof}
By the Runner-2 Elimination Lemma, runner~2 cannot absorb $I_9$, so runner~9 must first
close $I_4$. No alternative ordering is possible: any other runner covering $I_4$ in the
$\sigma_2=-1$ branch either falls outside the sign-range box (excluded by
Lemma~\ref{lem:sign-range}) or is already committed to a primary slot whose coverage
cannot be rerouted without opening a gap there. For runner~9 the $I_4$-closure condition is
\[
  \frac{43}{11a_9} = \frac{45}{110},\qquad a_9 = \frac{2(11k+10)}{9}.
\]
Inside $(9,10)$ the only possibility is $k=3$, giving $a_9=86/9$. The resulting interval
in $I_9$ is
\[
  B_{9,8}\!\left(\tfrac{86}{9}\right)
  = \left[\frac{441}{473},\,\frac{450}{473}\right].
\]
The complementary-pair cascade leaves four sign possibilities for $(\sigma_3,\sigma_4)$:

\medskip
\begin{center}
\begin{tabular}{cccc}
\toprule
$(\sigma_3,\sigma_4)$ & Uncovered region in $I_9$ & Width \\
\midrule
$(+,+)$ & $(741/770,\,272/275)$ & $103/3850$ \\
$(+,-)$ & $(10/11,\,441/473)$ together with the previous gap & $8279/165550$ \\
$(-,+)$ & $(450/473,\,1216/1265)$ & $538/54395$ \\
$(-,-)$ & $(10/11,\,441/473)$ together with $(450/473,\,1216/1265)$ & $1803/54395$ \\
\bottomrule
\end{tabular}
\end{center}

\medskip\noindent
All four widths are positive, and the minimum is $538/54395$. The detailed forced-speed
table and the endpoint comparisons are recorded in Appendix~A.11.
\end{proof}

\subsection{The branch \texorpdfstring{$\sigma_2=+1$}{sigma2 = +1}}

We now assume $\sigma_2=+1$. The pair $(2,9)$ carries the sharpest tension: closing $I_4$
drives $a_2$ upward, while bridging $I_9$ drives $a_2$ downward.

\begin{lemma}[The $I_4$ lower bound on $a_2$]\label{lem:i4lower}
Closing $I_4$ requires
\[
  a_2 \ge L(a_9) :=
  \begin{cases}
    20/9, & a_9 < 86/9,\\[2pt]
    2a_9/9, & a_9 \ge 86/9.
  \end{cases}
\]
(Note that $L$ is discontinuous at $a_9=86/9$: the left-hand limit is $20/9\approx 2.222$
while $L(86/9)=172/81\approx 2.123$. This downward jump is intentional: at $a_9=86/9$
runner~9 just reaches the left edge of $I_4$, reducing runner~2's required coverage.)
\end{lemma}

\begin{proof}
If $a_9<86/9$, runner~9 leaves the left edge $45/110$ uncovered in $I_4$.
Appendix~A.9 gives the full runner-by-runner exclusion table and shows that only runner~2
can cover $45/110$ without opening another slot, forcing
\[
  \frac{10}{11a_2} \le \frac{45}{110},\qquad a_2 \ge \frac{20}{9}.
\]
If $a_9\ge 86/9$, runner~9 reaches the left edge of $I_4$ but leaves a right subgap;
covering that subgap with runner~2 requires
\[
  \frac{10}{11a_2} \le \frac{45}{11a_9},\qquad a_2 \ge \frac{2a_9}{9}.
\]
\end{proof}

\begin{lemma}[The $I_9$ decomposition]\label{lem:i9decomp}
Under $\sigma_2=+1$, the parameter space splits into seven regions:

\medskip
\begin{center}
\begin{tabular}{llp{5cm}}
\toprule
Region & Condition & Consequence \\
\midrule
$\alpha_1$ & $\sigma_9=-1$, $a_9<580/69$ & direct incompatibility \\
$\alpha_2$ & $\sigma_9=-1$, $a_9\in[980/109,9)$ & direct incompatibility \\
$\alpha_3$ & $\sigma_9=+1$, $a_9\in(9,1960/207)$ & direct incompatibility \\
$\beta_1$  & $\sigma_9=-1$, $a_9\in[580/69,87/10)$ & right gap $[89/(11a_9),109/110]$ \\
$\beta_2$  & $\sigma_9=-1$, $a_9\in[87/10,980/109)$ & right gap $[23/(11a_2),109/110]$ \\
$\beta_3$  & $\sigma_9=+1$, $a_9\in[1960/207,86/9]$ & right gap $[100/(11a_9),109/110]$ \\
$\beta_4$  & $\sigma_9=+1$, $a_9\in(86/9,10)$ & right gap $[100/(11a_9),109/110]$ \\
\bottomrule
\end{tabular}
\end{center}
\end{lemma}

\noindent\textit{Exhaustiveness.} The seven regions partition the interval
$(a_9^{\min},a_9^{\max})=(8,10)$ (for $\sigma_9=-1$) and $(9,10)$ (for $\sigma_9=+1$)
exhaustively and disjointly: the $\alpha$-/$\beta$-boundary within each sign class is a
single threshold value, and adjacency of the subintervals is verified in Appendix~A.11.

\begin{proof}
We track two competing constraints on $a_2$.

\textit{Lower bound.} Lemma~\ref{lem:i4lower} gives $a_2\ge L(a_9)$.

\textit{Upper bound from $I_9$ coverage ($\sigma_9=-1$).} With $\sigma_9=-1$,
$a_9\in(8,9)$. Runner~9 uses $B_{9,7}=[87/(11a_9),89/(11a_9)]$ when $a_9<87/10$, and
$B_{9,8}=[98/(11a_9),100/(11a_9)]$ for $a_9\ge 980/109$. For runner~2's interval
$B_{2,1}=[23/(11a_2),25/(11a_2)]$ to abut $B_{9,7}$ without a gap:
\[
  \frac{23}{11a_2}\ge \frac{87}{11a_9} \implies a_2\le\frac{23a_9}{87}\quad
  (B_{9,7}\text{ regime}).
\]
Similarly, $a_2\le 23a_9/98$ in the $B_{9,8}$ regime and for $\sigma_9=+1$.

\textit{$\alpha$-regions: incompatibility.}
\begin{itemize}
  \item $\alpha_1$: $20/9>23a_9/87\iff a_9<580/69\approx 8.41$.
  \item $\alpha_2$: $20/9>23a_9/98\iff a_9<1960/207\approx 9.47$; throughout $[980/109,9)$.
  \item $\alpha_3$: same upper bound; incompatibility holds for $a_9\in(9,1960/207)$.
\end{itemize}

\textit{$\beta$-regions: residual right gap.} When the constraints are compatible,
the runners $B_{9,k}$ and $B_{2,1}$ together cover $I_9$ up to a residual right gap
with left endpoint $g_L\ge 100/110$ in each case. Hence Theorem~\ref{thm:ub} applies.
\end{proof}

\begin{theorem}[Positive-branch static core]\label{thm:poscore}
In every completely blocking configuration with $\sigma_2=+1$,
$\GF(I_k)>0$ for some slot $I_k$.
\end{theorem}

\begin{proof}
In the $\alpha$-regions of Lemma~\ref{lem:i9decomp}, the bounds from $I_4$ and $I_9$ are
incompatible. In the $\beta$-regions, Theorem~\ref{thm:ub} (Universal Blocking) shows that
the residual right gap either survives in $I_9$ or is absorbed only at the cost of opening
another slot via PSM, Chain~A, or Chain~B. In all cases $\GF(I_k)>0$ for some $k$.
\end{proof}

% ======================================================================
\section{Phase-Uniform Closure in \texorpdfstring{$\Wref$}{W\_ref}}
\label{sec:phase}
% ======================================================================

\subsection{Phase independence}

\begin{lemma}
For every runner $j$ and every cycle index $m$, $|B_{j,k}|=2/(11a_j)$.
Hence feasibility for covering a gap of fixed width depends only on $a_j$, not on phase.
\end{lemma}

\begin{proof}
Immediate from the formula for $B_{j,k}$.
\end{proof}

\begin{lemma}\label{lem:phaseindep2}
The hypotheses and conclusions of PSM, Chain~A, Chain~B, the Negative-Branch Gap Lemma,
and Theorem~\ref{thm:poscore} depend only on the speed ranges $a_j\in(j-1,j+1)$ and on
which runner covers the right gap of $I_9$. They are independent of the cycle index.
\end{lemma}

\begin{proof}
In pacemaker-relative coordinates the slot boundaries $I_s$ are fixed. In the $m$-th
pacemaker cycle, runner $j$'s relevant bad interval is $B_{j,k_m}$ where $k_m=k_0+m$.
The residual fractional position $\{a_jt - k_m\}$ equals the cycle-0 residual up to an
integer shift, so the slot-comparison inequalities are numerically identical in every cycle;
the $k$-index is invisible to the endpoint arithmetic. Every inequality in Sections~5
and~6 is a comparison involving only $a_j$ and the universal slot boundaries
$\{(11s+i)/110\}$. The branch choice (PSM~/ Chain~A~/ Chain~B) is determined by which
speed ranges hold and which runner covers the right gap of $I_9$; neither condition depends
on the absolute phase. No step depends on a specific cycle index $m$.
\end{proof}

\begin{remark}
The phase-independence of the $k$-index is not merely notational. It means that the
static-core analysis of Section~6 propagates unchanged to every complete pacemaker cycle
inside $\Wref$, and there is no hidden dependency on which $B_{j,k}$ index is serving as
the local primary branch in a given cycle.
\end{remark}

\subsection{The phase-uniform static core}

\begin{theorem}[Phase-uniform static core]\label{thm:phasecore}
Let $C_m\subset\Wref$ be any complete pacemaker cycle. Then some slot $I_k^{(m)}$ has
positive good fraction.
\end{theorem}

\begin{proof}
Split on $\sigma_2$. If $\sigma_2=+1$, Theorem~\ref{thm:poscore} gives $\GF(I_k)>0$ in
cycle-relative coordinates, and phase-independence makes this phase-uniform. If
$\sigma_2=-1$, the Negative-Branch Gap Lemma gives a positive witness width in $I_9$ by
a pure speed-range calculation, persistent in every cycle. Hence every complete pacemaker
cycle inside $\Wref$ contains a slot with positive good fraction.
\end{proof}

\begin{corollary}[$\Wref$ closure]\label{cor:wref}
In the commensurable large-ratio case $a_1<1$, there exists a witness in $\Wref$.
\end{corollary}

\begin{proof}
Runner~1 is isolated throughout $\Wref$ by definition. We need an integer $m$ such that
$C_m=[m/10,(m+1)/10]\subset\Wref$, i.e.\
\[
  \frac{10}{11a_1} \le m \le \frac{100}{11a_1}-1.
\]
The length of this interval is $90/(11a_1)-1>90/11-1=79/11>1$ because $a_1<1$. Hence it
contains an integer $m^*$, so $C_{m^*}\subset\Wref$. By Theorem~\ref{thm:phasecore},
some slot inside that cycle has positive good fraction for runners $2,\ldots,10$. Runner~1
is safe everywhere on the cycle, giving a witness for all runners.
\end{proof}

% ======================================================================
\section{The Irrational Case and the Full Proof}
% ======================================================================

\begin{lemma}[Commensurable reduction]\label{lem:comred}
If $q_1<\cdots<q_{10}$ are positive commensurable reals, then after clearing denominators
the problem reduces to the corresponding integer-speed configuration.
\end{lemma}

\begin{proof}
Write $q_i=p_i\omega$ with integers $p_i$ and $\omega>0$. Then
$\norm{q_it}=\norm{p_i(\omega t)}$,
so a witness for $(p_1,\ldots,p_{10})$ yields one for $(q_1,\ldots,q_{10})$.
\end{proof}

\noindent\textit{Logical note:} Corollary~\ref{cor:wref}, proved in \S\ref{sec:phase},
handles only the commensurable large-ratio case and is self-contained. There is no
circularity.

\begin{proof}[Proof of the Irrational-Ratio Proposition]
Let $\mathbf{x}(t)=(e_1t,\ldots,e_{10}t)\bmod 1\in\mathbb{T}^{10}$ and
$\mathcal{S}=(1/11,10/11)^{10}$. It suffices to show the orbit closure of $\mathbf{x}(t)$
meets $\mathcal{S}$.

Let $L=\{m\in\mathbb{Z}^{10}:m_1e_1+\cdots+m_{10}e_{10}=0\}$ and let $T_e$ be the
Kronecker subtorus determined by $L$. Since some ratio $e_i/e_j$ is irrational,
$\dim_{\mathbb{Q}}\operatorname{span}_{\mathbb{Q}}\{e_1,\ldots,e_{10}\}\ge 2$.

Choose a $\mathbb{Q}$-basis $b_1,\ldots,b_d$ from $\{e_1,\ldots,e_{10}\}$, write
$e_i=\sum_{k}r_{ik}b_k$ with $r_{ik}\in\mathbb{Q}$, pick a rational vector $v$ close to
$(b_1,\ldots,b_d)$, and define $q_i:=\sum_k r_{ik}v_k$. By taking $v$ sufficiently close
the $q_i$ are positive, distinct, ordered, and satisfy $q_{10}/q_1\ne 10$.

If $q_{10}/q_1<10$, the interval $(1/(11q_1),10/(11q_{10}))$ is nonempty; choose $s^*$
there. Then $y:=(\{q_1s^*\},\ldots,\{q_{10}s^*\})\in\mathcal{S}$.

If $q_{10}/q_1>10$, Lemma~\ref{lem:comred} and Corollary~\ref{cor:wref} apply to
$(q_1,\ldots,q_{10})$, producing an open witness set inside $\Wref$. Since $\mathcal{S}$
is open, any witness in that open set is strict; pick such a $y\in\mathcal{S}$.

In either case, the rational vector $v=(v_1,\ldots,v_d)$ was chosen within the rational
hyperplane defined by the integer relations $L$ of the original speed vector $(e_1,\ldots,e_{10})$:
explicitly, any integer relation $\sum c_i e_i=0$ (respectively $c_ie_i=c_je_j$) continues
to hold for the approximating speeds $q_i=\sum_k r_{ik}v_k$, because the coefficients
$r_{ik}\in\mathbb{Q}$ are unchanged and $v$ was perturbed within the same rational span.
Consequently the orbit closure of $(q_1,\ldots,q_{10})$ under the torus flow lies in $T_e$,
so $y\in T_e$.
By Kronecker's theorem~\cite{KN74}, the orbit $\{\mathbf{x}(t)\}$ is dense in $T_e$,
so some $t>0$ gives $\mathbf{x}(t)\in\mathcal{S}$.
\end{proof}

\begin{proof}[Proof of the Main Theorem]
Every speed configuration falls into exactly one of three cases:

\medskip
\begin{center}
\begin{tabular}{lll}
\toprule
Case & Condition & Method \\
\midrule
B   & $e_{10}/e_1\le 10$                         & Bounded-Ratio Theorem \\
E   & some $e_i/e_j\notin\mathbb{Q}$             & Irrational-Ratio Proposition \\
R$>$ & all $e_i/e_j\in\mathbb{Q}$ and $e_{10}/e_1>10$ & Corollary~\ref{cor:wref} \\
\bottomrule
\end{tabular}
\end{center}

\medskip\noindent
The three cases are mutually exclusive and exhaustive. Each is handled by the stated
method. Hence the Main Theorem holds for all ten-tuples of positive speeds.
\end{proof}

% ======================================================================
\begin{thebibliography}{99}

\bibitem{BGG+98}
A.~Bienia, L.~Goddyn, P.~Gvozdjak, A.~Sebo, M.~Tarsi,
\textit{Flows, view obstructions, and the lonely runner},
J.~Combin. Theory Ser.~B \textbf{72} (1998), 1--9.

\bibitem{BH01}
T.~Bohman, R.~Holzman,
\textit{Six lonely runners},
Electron.\ J.\ Combin.\ \textbf{8} (2001), R3.

\bibitem{BS08}
J.~Barajas, O.~Serra,
\textit{The lonely runner with seven runners},
Electron.\ J.\ Combin.\ \textbf{15} (2008), R48.

\bibitem{Cus73}
T.~W.~Cusick,
\textit{View obstruction problems},
Aequationes Math.\ \textbf{9} (1973), 165--170.

\bibitem{KN74}
L.~Kuipers, H.~Niederreiter,
\textit{Uniform Distribution of Sequences},
Wiley, 1974.

\bibitem{Pal26}
T.~Palelei,
\textit{The $\phi(n)$ law for arithmetic progressions in the lonely runner conjecture: algebraic structure and computational fragility},
Zenodo, 2026. \url{https://doi.org/10.5281/zenodo.18158886}

\bibitem{Ros25a}
B.~Rosenfeld,
\textit{The lonely runner conjecture holds for eight runners},
arXiv:2509.14111, 2025.

\bibitem{Ros25b}
B.~Rosenfeld,
\textit{The lonely runner conjecture holds for nine runners},
arXiv:2512.01912, 2025.

\bibitem{Tao18}
T.~Tao,
\textit{Some remarks on the lonely runner conjecture},
Contrib.\ Discrete Math.\ \textbf{13} (2018), 1--31.

\bibitem{Tra25}
T.~Trakulthongchai,
\textit{Nine and ten lonely runners},
arXiv:2511.22427, 2025.

\bibitem{Wil67}
J.~M.~Wills,
\textit{Zwei S\"atze \"uber inhomogene diophantische Approximation},
Monatsh.\ Math.\ \textbf{71} (1967), 263--269.

\end{thebibliography}

% ======================================================================
\appendix
% ======================================================================

\section{Technical Interval Arithmetic}

This appendix collects the endpoint arithmetic that is routine but too repetitive for the
main line. Sections~A.1--A.8 are the eight slot-boundary ledgers used in
Lemma~\ref{lem:sign-range}, with A.5--A.8 written as standalone right-edge calculations.
Section~A.9 records the extracted $45/110$ consequence used in Lemma~\ref{lem:i4lower}.
Section~A.10 records the feasibility table behind Universal Blocking.
Section~A.11 is the compact ledger for the Negative-Branch Gap Lemma.

\subsection*{A.1\quad Boundary ledger for $12/110$ (slot $I_1$ left)}

To straddle $12/110$, a branch $B_{j,k}$ must satisfy
\[
  \frac{11k+10}{11a_j} \le \frac{12}{110} \le \frac{11k+12}{11a_j}.
\]
Inside the sign-range box $a_j\in(j-1,j+1)$, the only off-schedule candidates are:

\begin{center}
\begin{tabular}{lll}
\toprule
branch $k$ & speed interval & candidates \\
\midrule
$0$ & $[25/3,\,10)$ & runners 8, 9 \\
\bottomrule
\end{tabular}
\end{center}

\begin{itemize}
  \item If runner~8 covers $12/110$, then $a_8\ge 25/3$ and
        $B_{8,2}[0]=32/(11a_8)\ge 32/99>34/110=I_3[\mathrm{left}]$,
        so the left edge of $I_3$ is missed.
  \item If runner~9 covers $12/110$, then either $a_9<86/9$ and
        $B_{9,3}[0]=43/(11a_9)>45/110$, or $a_9\ge 86/9$ and
        $B_{9,3}[1]=45/(11a_9)<54/110$, so $I_4$ is not closed.
\end{itemize}

Thus $12/110$ cannot be repaired off schedule without opening another claimed slot.

\subsection*{A.2\quad Boundary ledger for $23/110$ (slot $I_2$ left)}

\begin{center}
\begin{tabular}{lll}
\toprule
branch $k$ & speed interval & candidates \\
\midrule
$0$ & $[100/23,\,120/23]$ & runners 5, 6 \\
$1$ & $[210/23,\,10)$ & runner 9 \\
\bottomrule
\end{tabular}
\end{center}

\begin{itemize}
  \item Runners~5 or~6 via $B_{j,0}$: $B_{j,0}[0]\ge 23/132>12/110$, so $I_1$ acquires a
        left gap that neither runner can close alone.
  \item Runner~9 via $B_{9,1}$: same split as in A.1 shows runner~9 fails to close $I_4$.
\end{itemize}

\subsection*{A.3\quad Boundary ledger for $34/110$ (slot $I_3$ left)}

\begin{center}
\begin{tabular}{lll}
\toprule
branch $k$ & speed interval & candidates \\
\midrule
$0$ & $[50/17,\,3)$ & runner 2 \\
$0$ & $[3,\,60/17]$ & runner 4 \\
$1$ & $[105/17,\,115/17]$ & runners 6, 7 \\
$2$ & $[160/17,\,10)$ & runner 9 \\
\bottomrule
\end{tabular}
\end{center}

Each candidate either forces its pair-partner slot open or shifts runner~9's $I_4$ interval
off a required boundary. In each case one obtains positive good fraction in a partner slot.

\subsection*{A.4\quad Boundary ledger for $45/110$ (slot $I_4$ left)}

This is the calculation used in Lemma~\ref{lem:i4lower}.

\begin{center}
\begin{tabular}{lll}
\toprule
branch $k$ & speed interval & candidates \\
\midrule
$0$ & $[20/9,\,8/3]$ & runners 2, 3 \\
$1$ & $[14/3,\,46/9]$ & runners 4, 5, 6 \\
$2$ & $[64/9,\,68/9]$ & runners 7, 8 \\
$3$ & $[86/9,\,10]$ & runner 9 \\
\bottomrule
\end{tabular}
\end{center}

Runners~3--8 are each excluded: runner~3 forces $\GF(I_3)>0$; runner~4 triggers
Chain~B opening $I_7$; runner~5 leaves $I_1$ open; runner~6 leaves $I_8$ open;
runner~7 creates an inter-interval gap in $I_7$; runner~8 forces $\GF(I_3)>0$.
Runner~9 is excluded by the Lemma~\ref{lem:i4lower} left-gap analysis.
Therefore runner~2 is the unique viable filler, giving $a_2\ge 20/9$.

\subsection*{A.5\quad Boundary ledger for $65/110$ (slot $I_5$ right)}

\begin{center}
\begin{tabular}{lll}
\toprule
branch $k$ & speed interval & candidates \\
\midrule
$1$ & $[42/13,\,46/13]$ & runners 3, 4 \\
$2$ & $[64/13,\,68/13]$ & runners 4, 5, 6 \\
$3$ & $[86/13,\,90/13]$ & runners 6, 7 \\
$4$ & $[108/13,\,112/13]$ & runner 8 \\
\bottomrule
\end{tabular}
\end{center}

Every off-schedule repair of $65/110$ forces positive good fraction in a partner slot by
direct endpoint arithmetic (see the detailed exclusions in A.5 of the working notes).

\subsection*{A.6\quad Boundary ledger for $76/110$ (slot $I_6$ right)}

\begin{center}
\begin{tabular}{lll}
\toprule
branch $k$ & speed interval & candidates \\
\midrule
$0$ & $[25/19,\,30/19]$ & runner 2 \\
$1$ & $[105/38,\,3)$ & runner 2 \\
$1$ & $[3,\,115/38]$ & runner 4 \\
$2$ & $[80/19,\,85/19]$ & runners 4, 5 \\
$3$ & $[215/38,\,225/38]$ & runners 5, 6 \\
$4$ & $[135/19,\,140/19]$ & runner 7 \\
$5$ & $[325/38,\,335/38]$ & runner 9 \\
\bottomrule
\end{tabular}
\end{center}

Every off-schedule repair of $76/110$ leaves a partner slot open.

\subsection*{A.7\quad Boundary ledger for $87/110$ (slot $I_7$ right)}

\begin{center}
\begin{tabular}{lll}
\toprule
branch $k$ & speed interval & candidates \\
\midrule
$0$ & $[100/87,\,40/29]$ & runner 2 \\
$1$ & $[70/29,\,230/87]$ & runners 2, 3 \\
$2$ & $[320/87,\,340/87]$ & runner 3 \\
$3$ & $[430/87,\,150/29]$ & runners 5, 6 \\
$4$ & $[180/29,\,560/87]$ & runner 6 \\
$5$ & $[650/87,\,670/87]$ & runner 8 \\
$6$ & $[760/87,\,260/29]$ & runners 8, 9 \\
\bottomrule
\end{tabular}
\end{center}

Runner~2 ($k=0$): $B_{2,0}[\mathrm{left}]\ge 29/44>65/110$, so runner~2 misses $I_5$.
Runner~2 ($k=1$): $B_{2,1}[\mathrm{left}]\ge 1827/2530>56/110$, again missing $I_5$ left.
Runners~3 ($k=1,2$): the next interval $B_{3,k+1}$ jumps above $I_6$ right, and runner~8
alone cannot cover $I_6$. Runners~5 and~6 ($k=3$): the pair $(5,6)$ loses $I_8$ right
by incompatible bounds. Runner~6 ($k=4$), runner~8 ($k=5,6$), runner~9 ($k=6$): each
forces a partner slot open by analogous endpoint arithmetic.

\subsection*{A.8\quad Boundary ledger for $98/110$ (slot $I_8$ right)}

\begin{center}
\begin{tabular}{lll}
\toprule
branch $k$ & speed interval & candidates \\
\midrule
$0$ & $[50/49,\,60/49]$ & runner 2 \\
$1$ & $[15/7,\,115/49]$ & runners 2, 3 \\
$2$ & $[160/49,\,170/49]$ & runners 3, 4 \\
$3$ & $[215/49,\,225/49]$ & runner 4 \\
$5$ & $[325/49,\,335/49]$ & runner 7 \\
$6$ & $[380/49,\,390/49]$ & runners 7, 8 \\
$7$ & $[435/49,\,445/49]$ & runner 9 \\
\bottomrule
\end{tabular}
\end{center}

Runner~2 ($k=0$): $B_{2,0}[\mathrm{left}]\ge 49/66>67/110$, so $I_6$ left is uncovered.
Runners~2, 3 ($k=1$): the previous interval lies below $I_5$ or $I_6$, opening those slots.
Runner~4 ($k=2,3$): $B_{4,4}[\mathrm{left}]>98/110$ for $k=3$, so runner~4 misses $I_7$.
Runners~7, 8 ($k=6$): start above $I_6$ right, leaving $I_6$ left uncovered.
Runner~9 ($k=7$): $445/49<86/9$, so Lemma~\ref{lem:i4lower} forces $a_2\ge 20/9$,
but covering $56/110$ of $I_5$ requires $a_2\le 15/7<20/9$; contradiction.

\subsection*{A.9\quad The full $45/110$ exclusion used in Lemma~\ref{lem:i4lower}}

If $a_9<86/9$ and $I_4$ is to be closed, runner~9 leaves the left edge $45/110$ uncovered.
The A.4 exclusions show:
runner~3 opens $I_3$; runner~4 opens $I_7$ via Chain~B; runner~5 opens $I_1$;
runner~6 opens $I_8$; runner~7 opens $I_7$ through its inter-interval gap; runner~8 opens $I_3$.
Hence runner~2 is the unique viable filler and
\[
  \frac{10}{11a_2} \le \frac{45}{110} \implies a_2 \ge \frac{20}{9}.
\]

\subsection*{A.10\quad Feasibility table for Universal Blocking}

For a right gap $[g_L,109/110]\subset I_9$, the potentially feasible fillers are:

\begin{center}
\begin{tabular}{llll}
\toprule
Runner & Sign & Feasible range & Conclusion \\
\midrule
3 & $-1$ & $a_3\le 230/109$ & PSM opens $I_3$ \\
3 & $+1$ & $a_3\in[32/(11g_L),340/109]$ & Chain~A opens $I_8$ \\
4 & $-1$ & $a_4\le 340/109$ & PSM opens $I_2$ \\
4 & $+1$ & $a_4\in[43/(11g_L),450/109]$ & Chain~B opens $I_7$ \\
5 & $-1$ & $a_5\le 450/109$ & PSM opens $I_1$ \\
5 & $+1$ & no branch survives & no cover \\
6 & $\pm1$ & no branch survives & no cover \\
7, 8 & $\pm1$ & no branch survives & no cover \\
\bottomrule
\end{tabular}
\end{center}

More explicitly, for runners~5--8 with the positive sign (or any sign for runners~6--8),
the bounds
\[
  \frac{11k+10}{10} > \frac{10(11k+12)}{109}\quad\text{for all }k\ge 2
\]
rule out every feasible branch.

\subsection*{A.11\quad Compact ledger for the Negative-Branch Gap Lemma}

When $\sigma_2=-1$, the Runner-2 Elimination Lemma forces $a_9=86/9$ with
$B_{9,8}(86/9)=[441/473,450/473]$. The cascade-forced outer speeds by subcase:

\begin{itemize}
  \item $\sigma_3=+1$: $a_3=50/17$, $a_8=140/19$.
  \item $\sigma_3=-1$: $a_3=115/38$.
  \item $\sigma_4=+1$: $a_4=100/23$, $a_7=560/87$.
  \item $\sigma_4=-1$: $a_4=340/87$.
\end{itemize}

The four subcases are exactly the table in the Negative-Branch Gap Lemma, and the
minimal uncovered width is
\[
  \frac{1216}{1265} - \frac{450}{473} = \frac{538}{54395}.
\]

% ======================================================================
\section{Computational Corroboration}

The computations described here are \emph{not} part of the proof; the theorem is
established analytically in Sections~1--8.  They provide independent, large-scale
validation that the proof's structural mechanisms behave exactly as claimed, with no
hidden exceptional regimes detected.

\subsection*{Exact symbolic verification}

\begin{itemize}
  \item Exact rational checks on the canonical integer family $\{1,\ldots,9,p_{10}\}$
        with $p_{10}\in\{89,\ldots,200\}$ find witnesses in $\Wref$ in all 112 cases.
  \item All 34 explicit rational inequalities in the proof---including the endpoint
        constants
        \[
          \frac{83}{2695},\quad\frac{17}{430},\quad\frac{1}{50},\quad
          \frac{1533}{18490},\quad\frac{538}{54395}
        \]
        and the Chain~A, Chain~B, and negative-branch bounds---were verified with
        exact arithmetic in Python (\texttt{fractions.Fraction}), giving zero
        floating-point error.  All 34 checks pass.
\end{itemize}

\subsection*{Large-scale boundary-set verification}

Two independent witness-search campaigns were conducted on an Apple M4 MacBook Pro
(10-core CPU, Python~3.14).  The event-complete boundary-set method of~\cite{Pal26}
was used: rather than sampling a time-grid, the search evaluates only at the finite
set of times where some runner's fractional distance to the origin crosses the
$1/11$ threshold, guaranteeing that no witness in $\Wref$ can be missed between
evaluation points.  This approach, developed in the context of the $\phi(n)$ law for
arithmetic progressions~\cite{Pal26}, reduced per-configuration cost by a factor of
approximately $10^{3}$ compared with naive grid search and made the following scale
computationally feasible with no cloud resources.

\begin{center}
\begin{tabular}{lrrrrr}
\hline
Campaign & Seed & Configurations & Failures & Boundaries evaluated & Runtime \\
\hline
\texttt{run\_1e8\_mech} & 42   & $10^{8}$ & \textbf{0} & $6.1\times10^{10}$ & 65\,min \\
\texttt{run\_1e9\_mech} & 1337 & $10^{9}$ & \textbf{0} & $6.1\times10^{11}$ & 10.7\,hr \\
\hline
\textbf{Combined}       & ---  & $1.1\times10^{9}$ & \textbf{0} & $6.7\times10^{11}$ & --- \\
\hline
\end{tabular}
\end{center}

Each configuration was drawn randomly from the commensurable large-ratio regime
($a_1\in(0,1)$, $a_j$ uniformly in its sign-range interval for $j=2,\ldots,9$,
$a_{10}=10$) and classified into one of the eight analytic regions of the proof
($\sigma_2=\pm1$; sub-regions $\alpha_1,\alpha_3,\beta_1,\beta_2,\beta_3,\beta_4$).
Every region was covered in both campaigns and no exceptional behaviour was observed
in any region.

With zero failures across $1.1\times10^{9}$ event-complete trials, the 99.9\%
confidence upper bound on the true failure rate is $6.9\times10^{-9}$.  For
comparison, the probability of a hardware bit-flip error during the run exceeds this
bound, making undetected software error the dominant residual uncertainty.

All scripts, raw results, and reproduction instructions are available at:\\
\url{https://github.com/trentjp-gecta9/lrc-n11-computational-verification}

\end{document}
